\begin{abstract}
The growing adoption of autonomous large language model (LLM) agents with system and tool execution capabilities introduces a fundamental security challenge: enforcing fine-grained, intent-aware authorization while preserving low-latency execution. Existing approaches either rely on static permission models that do not adapt to dynamic execution contexts, or require synchronous semantic validation before each action, incurring substantial overhead. We present \SystemName, a runtime security framework for intent-driven authorization of computer-use agents. \SystemName introduces a Semantic-to-Kernel Intermediate Representation (SK-IR) that compiles high-level intents into verifiable low-level constraints enforceable via eBPF and Linux Security Modules (LSM), providing immediate kernel-level containment while semantic validation proceeds asynchronously. To reduce critical-path latency, \SystemName employs ShadowFS, a speculative execution layer that isolates file system and state modifications in an overlay, allowing operations to proceed before authorization completes. A trusted semantic guard evaluates intents in parallel, and changes are atomically committed upon approval or discarded upon rejection, mitigating time-of-check-to-time-of-use (TOCTOU) inconsistencies. \FIXME{We evaluate \SystemName across 15 categories of agent operations and show that it enforces least-privilege execution with low runtime overhead while significantly reducing perceived authorization latency through asynchronous validation.}
\end{abstract}